\section{Introduction}
\label{sec:intro}

The empirical literature on cartel detection has converged on
methods that read collusion off the distribution of submitted
bids: variance, kurtosis, skewness, or normalized spread within
a tender are informative about coordinated bidding when the
bid-level record is preserved
\citep{porter1993detection,bajari2003deciding,imhof2019detecting,wallimann2023machine}.
Real enforcement environments rarely match this informational
benchmark. Many electronic procurement systems archive only the
contract-award layer---winner identity, participant identity,
item code, and negotiated price---and export bid-by-bid records
to operational logs that competition authorities and audit
courts query slowly, partially, or not at all. The screening
methods designed for the bid layer become inoperative in the
layer that survives.

The question we ask is conceptual rather than operational: what
collusion-relevant information remains when enforcement observes
only the award layer? An answer cannot reproduce the bid-level
moment battery on coarsened data. It must instead identify a
statistic generated by the collusive arrangement itself that is
preserved when bid amounts are dropped. We argue that endogenous
loser-side participation is such a statistic. Cartels with cover
bidders generate persistent zero-win participation footprints
that are visible in award records, and the cumulative exposure
of those footprints carries information about cover-bidder type
that bid-distribution moments do not exhaust.

\paragraph{Theory}
A separating-equilibrium framework with cover bidders and
genuine entrants (Online Appendix~A) makes this concrete. Two
results anchor the empirical strategy. First, in equilibrium
cover bidders satisfy $\Pr(\text{win}\mid C)=0$: the
$\text{wins}=0$ filter is a separating choice rather than a
descriptive coincidence. Second, given award-record data, the
log-exposure statistic $\log(1+\text{tenders\_count})$ is a
sufficient ranking statistic for cover-bidder type under a
monotone-likelihood-ratio Poisson approximation; the binary
$\text{median}+1.5\!\times\!\text{IQR}$ rule is its
information-coarsening. The framework treats endogenous sorting
into cartel-affected environments as constitutive of the
screening object, not as a confound to be matched away. This
distinction governs the interpretation of every empirical
exercise that follows.

\paragraph{Empirical evidence}
We test the framework on São Paulo's Bolsa Eletr\^onica de
Compras (BEC, $\valYearStart$--$\valYearEnd$), an electronic
platform that preserves the award envelope for $\valSampleN$
tender-items across $1{,}308$ purchasing units and two procurement
modalities. Among the $\valBECfirms$ active firms,
$\valAlwaysLosers$ have a zero win rate across the full sample;
$\valFL$ of them cross the participation threshold the framework
identifies, and we call these \emph{frequent losers}. The
empirical content of the paper is structured by the screening
logic above. The continuous loss-intensity statistic strictly
dominates the binary rule in discrimination
(\citealt{delong1988comparing} test, $p<10^{-3}$), consistent
with the framework's identification of $\log(1+\text{tc})$ as
the primitive and the binary rule as its operational coarsening.
Frequent-loser presence associates with a $\valHeadlineRange$
higher conditional log negotiated unit price across four
estimators (cross-fit, OLS, CEM, IPW) on the broad sample. Once
the comparison is restricted to items where treated and control
overlap on observables the conditional coefficient reverses
sign: overlap-cell ATT $\valMatchOverlapCoef$,
propensity-score-trimmed $\valMatchPSCoef$. We read the reversal
as the empirical signature of a screening-value object rather
than a treatment-effect object: the overlap restriction strips
the endogenous sorting that, under the framework, generates the
screen's economic content. The broad-sample estimate is not a
defended causal effect; it is a measurement of the loser-side
participation footprint where the cartel deploys it.

\paragraph{Validation and scope}
Against CADE's procurement-cartel adjudication portfolio the
construct discriminates $\valCobidders$ adjudicated cobidder
firms inside the always-loser stratum at firm-level AUC
$\valAUCFLfirm$ in-sample, $\valAUCFLfirmTemp$ under temporal
holdout (train $2009$--$2016$, test $2017$--$2019$), and
$\valAUCprePost$ on a contemporaneous benchmark restricted to
four cases adjudicated before $2020$ ($95\%$ DeLong CI
$[0.713, 0.783]$). Cobidders are the structurally appropriate
validation object under coarsened observability: the
award-record envelope cannot adjudicate cartel membership, and
participation alongside adjudicated cartel members is the label
that the screening statistic is designed to recover. Against the
$\valDirectCADE$ direct CADE defendants active in the broader
BEC firm universe the construct discriminates at AUC
$\valAUCdirectStd$. The asymmetry between cobidder and
direct-defendant discrimination is the design's empirical
signature: direct defendants are typically frequent winners and
lie outside the construct's target stratum. The price imprint
varies across procuring-unit-size quartiles by a factor of
$\valPBUgradientRatio\!\!\times$, evidence of heterogeneity in
the detection regime that we interpret as scope information
rather than as identification of the institutional channel.

\paragraph{Contribution}
The contribution is at the conjunction of three literatures.
Bid-coordination tests \citep{bajari2003deciding} require a
researcher-supplied partition between colluders and a
competitive fringe; the construct supplies a candidate
\emph{ex ante} firm-level flag computable from award records.
Bid-distribution screens
\citep{imhof2019detecting,wallimann2023machine} require bid
microdata; the construct runs on the layer that survives the
operational coarsening. Empirical work on cover bidding has been
retrospective \citep{porter1999ohio,asker2010leniency}; the
construct shifts cover-bidder identification into a prospective
screening stage. The cumulative claim is structural: under
incomplete observability the appropriate enforcement architecture
separates an award-layer screening stage that produces candidate
flags from a bid-layer forensic stage that adjudicates suspected
coordination. The two stages answer different questions and
require different statistics; the construct fills the first.

\paragraph{What this paper does and does not do}
We do not claim a uniform treatment effect of cover bidding on
prices, identification of an oversight channel through which
buyer size moderates the participation footprint, or recovery of
cartel membership at the firm level. We claim that the
participation primitive identified by the framework carries
collusion-relevant information that survives the coarsening from
the bid layer to the award layer; that the screening-value
interpretation makes the sign reversal under overlap
restrictions diagnostic rather than confounding; and that the
construct discriminates a structurally appropriate
cartel-adjacency object at rates well above random matching at
comparable participation volume. The construct integrates with
bid-distribution methods rather than replacing them
(\S\ref{sec:rob_imhof}).

The remainder of the paper proceeds in five blocks.
Section~\ref{sec:literature} situates the contribution against
the bid-coordination, bid-distribution, and cover-bidding
literatures, and against work on enforcement design under
incomplete information. Section~\ref{sec:institutional}
describes the BEC data layers and the CADE adjudication
portfolio.
Sections~\ref{sec:data}--\ref{sec:empirical} state the
participation primitive, the binary operationalization, and the
descriptive identifying stance, alongside the two design-based
strategies (RDD, DiD) that return null.
Sections~\ref{sec:cade}--\ref{sec:results} report the CADE
validation and the within-item price association, with
\S\ref{sec:results_ols} treating the sign reversal under overlap
restrictions as the screening-value-vs-treatment-effect
distinction. Sections~\ref{sec:mechanisms}--\ref{sec:robustness}
develop the modal informativeness contrast, the buyer-size
heterogeneity, the leakage and adversarial-adaptation audits,
and the screening-vs-forensic-stage comparison against the
Imhof--Wallimann pipeline. Sections~\ref{sec:limitations}--\ref{sec:conclusion}
catalog scope limitations and conclude.
